% biography section
% 
% If you have an EPS/PDF photo (graphicx package needed) extra braces are
% needed around the contents of the optional argument to biography to prevent
% the LaTeX parser from getting confused when it sees the complicated
% \includegraphics command within an optional argument. (You could create
% your own custom macro containing the \includegraphics command to make things
% simpler here.)
%\begin{IEEEbiography}[{\includegraphics[width=1in,height=1.25in,clip,keepaspectratio]{mshell}}]{Michael Shell}
% or if you just want to reserve a space for a photo:
%\newpage
\vspace{-0.5 cm}



\begin{IEEEbiography}[{\includegraphics[width=1in,height=1.25in,clip,keepaspectratio]{photos/AlirezaMorsali.jpg}}]{AlirezaMorsali}
%\begin{IEEEbiographynophoto}{Luke Renaud}
received the B.Sc. and M.Sc. degrees in electrical engineering from University of Kerman, Iran, in 2009, and 2011, respectively. He finished his Ph.D. under supervision of Prof. Champagne in 2020 at McGill University, Montreal, Canada. His research interests include signal processing, machine learning and artificial intelligence. He was a recipient of several awards and scholarships, including, the McGill
Engineering Doctoral Award, STARaCom Collaborative Grant, and McGill BLUE Fellowship.
%\end{IEEEbiographynophoto}
\vspace{-0.8 cm}
\end{IEEEbiography}

%\vfill

%\begin{IEEEbiography}[{\includegraphics[width=1in,height=1.25in,clip,keepaspectratio]{photos/ylabmate}}]{Your Labmate}
\begin{IEEEbiographynophoto}{Your Labmate}
Bio paragraph
\end{IEEEbiographynophoto}
%\end{IEEEbiography}
%\vspace{-0.8 cm}

%\begin{IEEEbiography}[{\includegraphics[width=1in,height=1.25in,clip,keepaspectratio]{photos/yprof}}]{Your Professor}
\begin{IEEEbiographynophoto}{Your Professor}
Bio paragraph
\end{IEEEbiographynophoto}
%\end{IEEEbiography}

% insert where needed to balance the two columns on the last page with
% biographies
%\newpage

%\begin{IEEEbiographynophoto}{Jane Doe}
%Biography text here.
%\end{IEEEbiographynophoto}

% You can push biographies down or up by placing
% a \vfill before or after them. The appropriate
% use of \vfill depends on what kind of text is
% on the last page and whether or not the columns
% are being equalized.

\vfill

% Can be used to pull up biographies so that the bottom of the last one
% is flush with the other column.
%\enlargethispage{-5in}
